%% Master CV — Mahima Prajapati
%% Comprehensive overview of all competencies, experience, and achievements.
%% Use as the master reference when tailoring role-specific CVs.
%%
%% Compile with: pdflatex main_example.tex

\documentclass[11pt,a4paper]{moderncv}
\moderncvstyle{banking}
\moderncvcolor{blue}

\usepackage[utf8]{inputenc}
\usepackage{mathptmx}
\usepackage{hyperref}
\hypersetup{
    colorlinks=true,
    linkcolor=blue,
    filecolor=magenta,
    urlcolor=blue,
    pdftitle={Mahima Prajapati - CV},
    pdfpagemode=FullScreen,
}
\usepackage[scale=0.80]{geometry}
\usepackage{import}

% personal data
\name{Mahima}{Prajapati}
\address{Blacksburg, Virginia, USA}{}{}
\email{mahimap@vt.edu}
\extrainfo{\href{https://linkedin.com/in/mahimaprajapati}{LinkedIn}}

\begin{document}

\makecvtitle

% ============================================================
%     PROFILE STATEMENT
% ============================================================

\vspace{6pt}
\small{PhD candidate in Aerospace Engineering at Virginia Tech specializing in metal additive manufacturing process simulation and aeroelastic analysis. Expert in Abaqus FEA with hands-on experience predicting thermo-mechanical deformations during AM processes. Published AIAA SciTech researcher (2026); second paper in preparation for SciTech 2027. Proficient in MATLAB, SolidWorks, CATIA, and ANSYS, with internship experience in industrial automation and aircraft engine maintenance. Seeking Summer/Fall 2026 internship opportunities in aerospace structures, additive manufacturing, or FEA/simulation engineering.}

% ============================================================
%     CORE COMPETENCIES
% ============================================================

\section{Core Competencies}
\vspace{1pt}
\begin{itemize}

\item \textbf{Finite Element Analysis (FEA):} Abaqus (expert) — thermo-mechanical coupled simulation, metal additive manufacturing process modeling, deformation and distortion prediction, UMAT/VUMAT subroutines; ANSYS — structural analysis and validation.
\vspace{1pt}

\item \textbf{CAD / CAM:} SolidWorks — 3D mechanical design and assembly; CATIA — aerospace-grade parametric design; AutoCAD Electrical — electrical systems documentation.
\vspace{1pt}

\item \textbf{Computational Analysis:} MATLAB and Simulink — numerical methods, data processing, dynamic system modeling; scripting for simulation workflows.
\vspace{1pt}

\item \textbf{Aerospace Structures \& Aeroelasticity:} Flutter analysis, fluid-structure interaction, structural dynamics, vibration studies; MS thesis on aeroelastic instability of wings with distributed electric propulsors.
\vspace{1pt}

\item \textbf{Technical Communication:} AIAA-published researcher; experienced instructor and technical mentor to 50+ engineering students; LaTeX documentation for academic and professional materials.

\end{itemize}

% ============================================================
%     PROFESSIONAL EXPERIENCE
% ============================================================

\section{Professional Experience}
\vspace{3pt}
\begin{itemize}

% --- Most Recent Role ---
\item{\cventry{Dec 2021--Present}{Graduate Research \& Teaching Assistant}{Virginia Tech, Kevin T. Crofton Dept. of Aerospace \& Ocean Engineering}{Blacksburg, VA}{}{\vspace{1pt}
\begin{itemize}
    \item Developing Abaqus-based thermo-mechanical FEA simulations to predict deformations and distortions during metal additive manufacturing processes; bridging computational modeling with real-world manufacturing outcomes.
    \vspace{1pt}
    \item Conducted MS research on aeroelastic flutter speed and frequency of aircraft wings with distributed electric propulsors; published and presented findings at AIAA SciTech 2026.
    \vspace{1pt}
    \item Delivered lectures, tutorials, and laboratory sessions for undergraduate aerospace engineering courses; provided hands-on technical guidance to 50+ students per semester.
    \vspace{1pt}
    \item Contributed to vibrational studies of the Broad Art Museum Monumental Staircase System with Dr.\ Mehdi Setareh (2021).
\end{itemize}}}

\vspace{3pt}

% --- Automation Internship ---
\item{\cventry{Jun 2019--Jul 2019}{Mechanical Engineering Intern}{Automation \& Control Systems (now PLC Automation Training)}{Pune, India}{}{\vspace{1pt}
\begin{itemize}
    \item Improved efficiency and accuracy of a 6-axis ABB Robot by 20\% and reduced downtime by 15\% using RobotStudio 6.00; performed circuit error diagnosis with CoDeSys and RSlogix 8.10.00.
    \vspace{1pt}
    \item Redesigned complex electrical projects using AutoCAD Electrical, saving 70\% of working cost.
\end{itemize}}}

\vspace{3pt}

% --- Air India Internship ---
\item{\cventry{Dec 2018--Jan 2019}{Engineering Intern}{Air India Limited}{Mumbai, India}{}{\vspace{1pt}
\begin{itemize}
    \item Assisted on-site engineers in routine maintenance of aircraft engines (PW4056, CFM56, GE90/GE~Nx) and APU at Mumbai International Airport; contributed to 20\% efficiency improvement of PW4056 engine.
    \vspace{1pt}
    \item Gained hands-on experience with engine assembly modules across commercial jet platforms.
\end{itemize}}}

\end{itemize}

% ============================================================
%     EDUCATION
% ============================================================

\section{Education}
\vspace{1pt}
\begin{itemize}

\item{\cventry{Jan 2025--Present}{Doctor of Philosophy --- Aerospace \& Ocean Engineering}{Virginia Tech}{Blacksburg, VA}{}{\vspace{1pt}
Research: Metal additive manufacturing process simulation; Abaqus-based FEA modeling of thermo-mechanical deformations and distortions during AM.
}}

\vspace{3pt}

\item{\cventry{Aug 2021--Dec 2024}{Master of Science --- Aerospace Engineering}{Virginia Tech}{Blacksburg, VA}{}{\vspace{1pt}
Thesis: ``Aeroelastic Analysis of a Uniform Cantilever Wing with Distributed Electric Propulsors with Rotary Inertia.'' Published at AIAA SciTech 2026 (\href{https://doi.org/10.2514/6.2026-1649}{DOI: 10.2514/6.2026-1649}).
}}

\vspace{3pt}

\item{\cventry{2016--2020}{Bachelor of Engineering --- Mechanical Engineering}{M.H.\ Saboo Siddik College of Engineering}{Mumbai, India}{}{\vspace{1pt}
}}

\end{itemize}

% ============================================================
%     LANGUAGES
% ============================================================

\section{Languages}
\vspace{1pt}
\begin{itemize}
\item English (professional), Hindi (professional), Gujarati (full professional), Marathi (limited working), Spanish (elementary).
\end{itemize}

% ============================================================
%     PUBLICATIONS
% ============================================================

\section{Publications}
\vspace{1pt}
\begin{itemize}
\item Prajapati, M.~D., and Kapania, R.~K. (2026). ``Aeroelastic Analysis of a Uniform Cantilever Wing with Distributed Electric Propulsors with Rotary Inertia.'' \textit{AIAA SciTech Forum 2026}. \href{https://doi.org/10.2514/6.2026-1649}{DOI: 10.2514/6.2026-1649}
\item Prajapati, M.~D., and Kapania, R.~K. (2027). ``Effect of Thrust on Aeroelastic Analysis of DEP Wings with Rotary Inertia.'' \textit{AIAA SciTech Forum 2027}. (In preparation; journal paper forthcoming.)
\end{itemize}

% ============================================================
%     HONORS AND AWARDS
% ============================================================

\section{Honors and Awards}
\vspace{1pt}
\begin{itemize}
\item{\textbf{First Place --- Innovation \& Technology Advancement} -- Virginia Tech Graduate \& Professional Student Symposium (GPSS); competitive poster award for aeroelastic flutter research; \$275 monetary award.}
\item{\textbf{Nutshell Games 2024} -- Virginia Tech Center for Communicating Science; presented ``Wings in Motion: Mass and Flutter Dynamics'' to a public panel in 90 seconds.}
\item{\textbf{Future Professoriate Certificate} -- Virginia Tech Graduate School.}
\item{\textbf{Tolson \& Nerney Fellowship} -- Kevin T. Crofton Dept.\ of Aerospace \& Ocean Engineering, Virginia Tech (2025); \$11,750.}
\end{itemize}

% ============================================================
%     REFERENCES
% ============================================================

\section{References}
\vspace{1pt}
\begin{itemize}
\item{Available upon request.}
\end{itemize}

\end{document}
